\documentclass[a4paper]{article}
\usepackage{amsfonts}
\usepackage{amsmath}
\usepackage{amssymb}
\usepackage{graphicx}    

\begin{document}
  
 
\noindent \large Auteur: Elaine de Bie\\
\noindent \large Studierichting: Informatica\\
\noindent \large Oefeningenreeks 1E nr.16 \\

\medskip


\normalsize


Als $i$ en $2i$ nulpunten zijn, dan volgen daaruit ook de nulpunten $-i$ en $-2i$.\\

$A(z) = a\cdot(z-i)\cdot(z+i)\cdot(z-2i)\cdot(z+2i)$ met $a \neq 0$\\

$=a(z^2 + iz - iz +1)\cdot(z^2 +2iz - 2iz +4)$\\

$=a(z^2 + 1)\cdot(z^2 + 4)$\\

$=a [(z^2 + 1)(z^2) + (z^2 + 1)(4)]$\\

$=a (z^4 + z^2 + 4z^2 + 4)$\\

$=a (z^4 + 5z^2 + 4)$\\

$=az^4 + 5az^2 + 4a$\\



\begin{thebibliography}{999}
%\bibitem{a} Referentie
\end{thebibliography}
\end{document}
